% Informe de la Fase 2. Edición final e integración: Persona 1.
% Cada sección está en secciones/ para que cada persona escriba la suya sin conflictos de Git.
%
% Compilar:  latexmk -pdf informe.tex      (o pdflatex dos veces)
\documentclass[11pt, a4paper]{article}

\usepackage[T1]{fontenc}
\usepackage[spanish, es-noshorthands, es-tabla]{babel}
\usepackage{lmodern}
\usepackage{microtype}
\usepackage{upquote}
\usepackage[margin=2.5cm]{geometry}
\usepackage{amsmath, amssymb, amsthm}
\usepackage{booktabs}
\usepackage{graphicx}
\usepackage{xcolor}
\usepackage{listings}
\usepackage{tikz}
\usetikzlibrary{automata, arrows.meta, positioning}
\usepackage[hidelinks]{hyperref}

\graphicspath{{../../results/figures/}{../img/}}

\lstset{basicstyle=\ttfamily\small, breaklines=true, columns=fullflexible,
        literate={á}{{\'a}}1 {é}{{\'e}}1 {í}{{\'i}}1 {ó}{{\'o}}1 {ú}{{\'u}}1 {ñ}{{\~n}}1}

% Marca visible para lo pendiente; quitar antes de entregar.
\newcommand{\pendiente}[1]{\textcolor{red}{\textbf{[PENDIENTE: #1]}}}
% Árbol entre paréntesis en texto.
\newcommand{\arbol}[1]{\texttt{#1}}

\title{Fase 2: análisis sintáctico de un fragmento del español\\
  con una gramática libre de contexto}
\author{%
  Pablo Vides (00080323) \and Gabriela Valiente (00010122) \and
  Valeria Ortiz (00206022) \and Paola Pérez Pérez (00117222) \and
  Alejandro Melgar (00026222)}
\date{\pendiente{fecha de entrega}}

\begin{document}
\maketitle

\begin{abstract}
\pendiente{P1: resumen de 150--200 palabras al terminar el informe.}
\end{abstract}

\tableofcontents
\newpage

% Responsable: P1
\section{Marco teórico}
\label{sec:marco_teorico}

\pendiente{P1}

% Contenido esperado:
\begin{itemize}
  \item Gramáticas libres de contexto: definición formal $G = (V, \Sigma, R, S)$, derivaciones y árboles.
  \item Jerarquía de Chomsky y dónde queda el lenguaje natural; enlace con la reflexión de la Fase 1 (\{wcw\}).
  \item Rasgos y concordancia: por qué multiplicar categorías hace crecer la gramática.
  \item Análisis descendente recursivo y retroceso (backtracking).
  \item Enfoques estadísticos y neuronales (spaCy, Transformers, LLM): qué aprenden y qué no garantizan.
\end{itemize}

% Responsable: P1
\section{Diseño de la gramática}
\label{sec:gramatica}

\pendiente{P1}

% Contenido esperado:
\begin{itemize}
  \item Producciones, categorías y rasgos (desde docs/gramatica\_es.md).
  \item Cobertura: SVO, adjetivos, adverbios, frases preposicionales, dos tiempos, concordancia.
  \item Decisiones de diseño y ambigüedades conocidas.
  \item Análisis formal de fallas: qué no puede describir la gramática y por qué.
\end{itemize}

% Responsable: P2  (versión completa en extendido/03_parser.tex)
\section{Implementación del parser}
\label{sec:parser}

El parser (\texttt{spanish/parser.py}) es descendente recursivo: cada no terminal es un método
(\texttt{sentence}, \texttt{noun\_phrase}, \texttt{verb\_phrase}…). Las reglas con
$\varepsilon$ son un \texttt{if} y las reglas $X' \to \alpha X' \mid \varepsilon$ son un
\texttt{while}, ambos guiados por el conjunto FIRST de la palabra actual. Una oración se acepta
solo si se construye el árbol y no sobran palabras.

\paragraph{Palabras ambiguas.} El lexer entrega todas las lecturas de cada palabra y el parser
toma la que pide su posición: en \emph{El gato duerme bajo la mesa}, \emph{bajo} es preposición;
en \emph{El hombre bajo come}, adjetivo. Sin retroceso, esto falla en un caso que la gramática sí
genera: en \emph{El gato bajo la mesa duerme} el parser toma \emph{bajo} como adjetivo y luego no
puede seguir. Resolverlo requiere retroceso o un parser de gráfico.

\paragraph{Pipeline.} \texttt{spanish/pipeline.py} encadena lexer, parser, analizador de
concordancia y visualización; si una etapa falla, la siguiente no se ejecuta. Cada error indica
la palabra y la causa:
\begin{center}
  \small
  \resizebox{\linewidth}{!}{%
  \begin{tabular}{ll}
    \toprule
    Oración & Salida \\
    \midrule
    \emph{El perro corre en el parque.} & sintáctico, palabra 3: `corre' no está en la gramática \\
    \emph{Come la manzana.} & sintáctico, palabra 1: se esperaba un pronombre, determinante o posesivo \\
    \emph{Los perros come.} & concordancia, palabra 3: el verbo `come' (singular) no concuerda con `perros' (plural) \\
    \bottomrule
  \end{tabular}}
\end{center}
El resultado incluye además las etiquetas, el sujeto, el verbo, el objeto y la unión de la
frase preposicional, que se sacan del árbol. Hay 220 pruebas automáticas; una de ellas pasa las
73 oraciones del conjunto de prueba por el pipeline completo.

% Responsable: P3  (versión completa en extendido/04_lexer_concordancia_arboles.tex)
\section{Léxico, concordancia y visualización}
\label{sec:lexer_concordancia_arboles}

\paragraph{Lexer.} Es el de la Fase~1 adaptado: el alfabeto incluye \emph{á, é, í, ó, ú, ü, ñ},
el texto se normaliza a Unicode NFC, y la categoría ya no sale de una expresión regular sino del
léxico, porque la misma forma puede ser sustantivo, verbo o adjetivo. Una palabra desconocida
no es un error léxico: recibe la categoría \texttt{X} y el parser rechaza la oración.

\paragraph{Léxico.} \texttt{data/lexicon.json} tiene 221 palabras con su categoría y sus rasgos
(género, número, persona y tiempo): 52 sustantivos, 58 adjetivos, 7 verbos en presente y
pretérito (70 formas) y las palabras funcionales. Cinco palabras son ambiguas, como \emph{bajo}.

\paragraph{Analizador semántico (concordancia).} \texttt{spanish/agreement.py} es la versión de
la Fase~2 del analizador semántico (\texttt{semantico.py}) y la tercera etapa del pipeline. Así
como en la Fase~1 la tabla de símbolos verificaba reglas que la CFG no expresa, este analizador
recorre el árbol y verifica que, en cada sintagma nominal, determinante, sustantivo y adjetivo coincidan en género y número, y que
el núcleo del sujeto coincida con el verbo en número y persona. Un rasgo ausente concuerda con
todo (adjetivos invariables como \emph{grande}, posesivos sin género). Como trabaja sobre el
árbol, resuelve la concordancia a distancia: acepta \emph{El perro de los niños come} y rechaza
\emph{*El perro de los niños comen}.

\paragraph{Visualización.} \texttt{spanish/trees.py} dibuja el árbol en la terminal
(\texttt{nltk}) o como imagen con Graphviz, con un color por categoría y las palabras que no
concuerdan en rojo
(figura~\ref{fig:arbol_concordancia}).

\begin{figure}[h]
  \centering
  \includegraphics[width=0.42\linewidth]{arbol_concordancia}
  \caption{Árbol generado por el pipeline para \emph{*La elefante grises duerme}: se deriva,
    pero el analizador semántico marca en rojo las dos palabras que no concuerdan.}
  \label{fig:arbol_concordancia}
\end{figure}

% Responsable: P4  (versión completa en extendido/05_set_de_prueba.tex)
\section{Conjunto de prueba y metodología}
\label{sec:set_de_prueba}

El conjunto (\texttt{data/test\_set.csv}) tiene 73 oraciones en tres niveles
(tabla~\ref{tab:set_de_prueba}). Cada una lleva la salida esperada: si la gramática debe
aceptarla, si es español correcto, el tipo de error, los fenómenos que prueba, la categoría de
cada palabra, el sujeto, el verbo, el objeto y a qué se une la primera frase preposicional.
\begin{itemize}
  \item \textbf{Nivel 1 (simples):} sujeto-verbo y sujeto-verbo-objeto, todas aceptables.
  \item \textbf{Nivel 2 (intermedias):} adjetivos, intensificadores, adverbios y frases
    preposicionales, más errores de concordancia (\emph{*La perro come}).
  \item \textbf{Nivel 3 (complejas):} varias frases preposicionales, ambigüedad estructural,
    concordancia a distancia y 11 oraciones correctas que la gramática no genera.
\end{itemize}

\begin{table}[h]
  \centering
  \small
  \resizebox{\linewidth}{!}{%
  \begin{tabular}{lrrrrrr}
    \toprule
    Nivel & Oraciones & En la gramática & Válidas fuera & Error sintáctico & Error de concordancia & Con PP \\
    \midrule
    1 & 20 & 20 & 0 & 0 & 0 & 0 \\
    2 & 29 & 18 & 0 & 1 & 10 & 6 \\
    3 & 24 & 7 & 11 & 14 & 3 & 11 \\
    Total & 73 & 45 & 11 & 15 & 13 & 17 \\
    \bottomrule
  \end{tabular}}
  \caption{Composición del conjunto de prueba.}
  \label{tab:set_de_prueba}
\end{table}

\paragraph{Dos referencias.} Cada decisión se evalúa contra \texttt{en\_gramatica} (¿sigue las
reglas?) y contra \texttt{espanol\_valido} (¿es español correcto?). Las 11 oraciones correctas
fuera de la gramática separan a un sistema que sigue las reglas de uno que sigue su conocimiento
del español.

\paragraph{Métricas.} Con la clase positiva \emph{acepta}:
precisión $= VP/(VP+FP)$, cobertura $= VP/(VP+FN)$, rechazo correcto $= VN/(VN+FP)$ y
exactitud. Además: tipo de error correcto en los rechazos, exactitud del etiquetado palabra por
palabra, exactitud de sujeto, verbo y objeto, y exactitud de la unión de la frase preposicional
(una falta de respuesta cuenta como error). Cada sistema participa solo donde produce datos.

\paragraph{Sistemas.} Parser propio; spaCy 3.8 (\texttt{es\_core\_news\_sm}); Transformers
(XLM-R para etiquetar y BETO para la gramaticalidad); y el LLM \texttt{gemini-3.5-flash-lite}
con temperatura~0, en versiones \emph{zero-shot} y \emph{few-shot}. Todos se corrieron el 26 de
noviembre de 2026; las métricas salen de \texttt{python -m evaluation.metrics}.

% Responsable: P4
% Tablas y figuras: python -m evaluation.metrics  (results/tablas/*.tex, results/figures/*.png).
% \tablagenerada está definida en 05_set_de_prueba.tex.
\section{Resultados}
\label{sec:resultados}

\newcommand{\figuragenerada}[2][\linewidth]{%
  \IfFileExists{../../results/figures/#2.png}{\includegraphics[width=#1]{#2}}{%
    \pendiente{generar results/figures/#2.png}}}

\pendiente{P4: estas tablas solo tienen el parser propio. Falta volver a correr spaCy y
Transformers (P1) con el conjunto de 73 oraciones y el LLM (P5); después se regeneran con
python -m evaluation.metrics y se completa la discusión.}

\subsection{Aceptación}

\begin{table}[h]
  \centering
  \small
  \tablagenerada{aceptacion}\par\medskip
  \tablagenerada{aceptacion_espanol}
  \caption{Aceptación contra \texttt{en\_gramatica} (arriba) y contra \texttt{espanol\_valido}
    (abajo). ``---'': la métrica no aplica al sistema.}
  \label{tab:aceptacion}
\end{table}

\begin{figure}[h]
  \centering
  \figuragenerada{aceptacion}
  \caption{Métricas de aceptación contra las reglas de la gramática.}
  \label{fig:aceptacion}
\end{figure}

El parser propio acepta exactamente las 45 oraciones de la gramática y rechaza las 28 que no
lo están. Además, todas las rechaza por la razón esperada: 15 por error sintáctico y 13 por
concordancia (tabla~\ref{tab:tipo_error}). Este resultado no mide qué tan bueno es el enfoque,
porque el conjunto se etiquetó con la misma gramática que implementa el parser. Lo que muestra
es que la implementación es fiel a la especificación, incluso en los casos difíciles: la
concordancia a distancia (\emph{El perro de los niños come} se acepta y \emph{*...comen} se
rechaza) y la palabra ambigua \emph{bajo}.

La comparación interesante es contra \texttt{espanol\_valido}. Ahí el parser tiene precisión
1.00, pero su cobertura baja a 0.80. Nunca acepta español incorrecto, pero rechaza 11 de las 56
oraciones correctas: todo lo que la CFG no genera (sujeto omitido, orden libre, coordinación,
subordinadas y palabras fuera del léxico). Ese 20\,\% es la medida concreta de cuánto español
correcto queda fuera de la gramática.

\begin{figure}[h]
  \centering
  \figuragenerada[0.7\linewidth]{reglas_vs_espanol}
  \caption{Exactitud de la aceptación contra las dos referencias.}
  \label{fig:reglas_vs_espanol}
\end{figure}

\pendiente{P4 con P5: comparar con Transformers y el LLM. Si su exactitud contra
espanol\_valido supera a la de en\_gramatica, siguen su conocimiento del español y no las reglas.}

\begin{table}[h]
  \centering
  \small
  \resizebox{\linewidth}{!}{\tablagenerada{tipo_error}}
  \caption{Tipo de error esperado frente a la respuesta del sistema, en las 28 oraciones fuera
    de la gramática.}
  \label{tab:tipo_error}
\end{table}

\subsection{Resultados por nivel}

\begin{table}[h]
  \centering
  \small
  \tablagenerada{exactitud_aceptacion_por_nivel}\par\medskip
  \tablagenerada{etiquetado_por_nivel}
  \caption{Exactitud de la aceptación (arriba) y del etiquetado (abajo) por nivel.}
  \label{tab:por_nivel}
\end{table}

\begin{figure}[h]
  \centering
  \figuragenerada{por_nivel}
  \caption{Exactitud por nivel de complejidad.}
  \label{fig:por_nivel}
\end{figure}

\pendiente{P4: describir cómo cambia cada sistema del nivel 1 al 3.}

\subsection{Etiquetado y funciones sintácticas}

\begin{table}[h]
  \centering
  \small
  \tablagenerada{etiquetado}\quad\tablagenerada{funciones}
  \caption{Exactitud del etiquetado (y oraciones con otro número de palabras), y de las
    funciones sintácticas.}
  \label{tab:etiquetado}
\end{table}

\begin{table}[h]
  \centering
  \small
  \tablagenerada{etiquetado_categoria}
  \caption{Cobertura (\emph{recall}) del etiquetado por categoría.}
  \label{tab:etiquetado_categoria}
\end{table}

\begin{figure}[h]
  \centering
  \figuragenerada[0.8\linewidth]{etiquetado_categoria}
  \caption{Cobertura del etiquetado por categoría.}
  \label{fig:etiquetado_categoria}
\end{figure}

El parser etiqueta sin errores, y da bien el sujeto, el verbo y el objeto, en todas las
oraciones que puede derivar (58 de 73: las aceptadas y las que fallan solo en la
concordancia). En las otras 15 no produce etiquetas. Un etiquetador estadístico sí puede
analizarlas, y esa diferencia de cobertura es la que se discute en la
sección~\ref{sec:comparacion_nlp}.

\subsection{Unión de la frase preposicional}

\begin{table}[h]
  \centering
  \small
  \tablagenerada{pp_adjunta}\par\medskip
  \tablagenerada{pp_confusion}
  \caption{Unión de la primera frase preposicional en las 17 oraciones que la tienen.}
  \label{tab:pp}
\end{table}

\begin{figure}[h]
  \centering
  \figuragenerada{funciones_pp}
  \caption{Funciones sintácticas y unión de la frase preposicional.}
  \label{fig:funciones_pp}
\end{figure}

Aquí aparece el costo del conflicto LL(1) de la sección~\ref{sec:gramatica}. El parser une la
frase preposicional al sustantivo más cercano siempre que hay un objeto. Por eso acierta en
las 6 oraciones donde la unión correcta es al sustantivo, y falla en las 4 donde un hablante la
une al verbo (\emph{Ellos ven los árboles altos desde la casa}). Las 2 oraciones ambiguas las
resuelve sin avisar, como si no lo fueran, así que detecta 0 de 2 ambigüedades. Su exactitud es
0.59 (10 de 17). Una de las fallas es \emph{El perro corre en el parque}, que no se puede
derivar porque \emph{corre} no está en el léxico.

\subsection{Costo}

\begin{table}[h]
  \centering
  \small
  \tablagenerada{costo}
  \caption{Tiempo medio por oración y fallos de ejecución.}
  \label{tab:costo}
\end{table}

\pendiente{P4: comparar el tiempo del parser (centésimas de milisegundo por oración) con el de
spaCy, Transformers y el LLM cuando estén sus resultados.}

% Responsable: P1
\section{Comparación con herramientas de NLP}
\label{sec:comparacion_nlp}

\pendiente{P1}

% Contenido esperado:
\begin{itemize}
  \item Qué analizan bien spaCy y Transformers que la CFG no puede (y por qué).
  \item Qué acierta la CFG y ellos no (concordancia).
  \item Ejemplos concretos del notebook nlp\_comparison.ipynb.
\end{itemize}

% Responsable: P5  (versión completa en extendido/08_analisis_llm.tex)
\section{Experimento con el LLM}
\label{sec:analisis_llm}

Le dimos a \texttt{gemini-3.5-flash-lite} (temperatura 0) la gramática sin recursión por la
izquierda, el léxico, las formas de árbol permitidas y la instrucción de usar \emph{solo} esas
reglas ``aunque conozcas otras construcciones válidas del español''; debía responder un JSON
con la decisión, el tipo de error, las etiquetas, las funciones y el árbol. La versión
\emph{few-shot} agrega cinco ejemplos resueltos que no están en el conjunto de prueba. Los 146
prompts y respuestas están en \texttt{results/logs/llm.jsonl}.

\begin{table}[h]
  \centering
  \small
  \begin{tabular}{lrr}
    \toprule
    & \emph{Zero-shot} & \emph{Few-shot} \\
    \midrule
    Exactitud contra las reglas / contra el español & 0.97 / 0.82 & 0.99 / 0.86 \\
    Acepta español correcto fuera de la gramática & 0 / 11 & 1 / 11 \\
    Acepta español incorrecto & 0 / 17 & 0 / 17 \\
    Árboles con las formas de la gramática & 67\,\% & 96\,\% \\
    Árboles iguales al del parser & 49\,\% & 82\,\% \\
    \bottomrule
  \end{tabular}
  \caption{Resultados del LLM; los árboles se evalúan en las oraciones que acepta.}
  \label{tab:llm}
\end{table}

\paragraph{(a) ¿Respeta las reglas?} Sí, en la decisión. Su exactitud contra las reglas supera a
la del español, al contrario de Transformers: rechaza \emph{El perro y el gato duermen} o
\emph{La niña que lee un libro come}, que sabe que son español correcto. Pero cuando fuerza una
regla puede aplicarla mal. La única desviación del \emph{few-shot} acepta \emph{Leímos un
libro} con el árbol \arbol{(S (NP un libro) (VP leímos))}: convierte el objeto en sujeto para
encajar en $S \to NP\ VP$, sin notar que no concuerdan.

\paragraph{(b) ¿Dónde acierta y dónde falla?} Decide bien casi todo el nivel 1 y 2; sus errores
están en oraciones con \emph{sus pijamas}, que \emph{zero-shot} rechaza aunque son correctas, y
en el tipo de error de \emph{Come la manzana}. Es el mejor en etiquetado (0.99) y en la unión de
la frase preposicional (15 de 17), pero nunca usa la respuesta \texttt{ambigua}, aunque el prompt
la ofrece. Sus árboles son el punto débil: en \emph{zero-shot}, uno de cada tres no respeta la
gramática (por ejemplo, \arbol{(NPbase ella)} sin la categoría \textit{Prn}).

\paragraph{(c) ¿Mejora el \emph{few-shot}?} Sobre todo en la forma: los árboles válidos pasan del
67\,\% al 96\,\%. La decisión mejora poco (de 0.97 a 0.99) y aparece la desviación de
\emph{Leímos un libro}.

% Responsable: P5
\section{Comparación de enfoques}
\label{sec:comparacion_enfoques}

\pendiente{P5}

% Contenido esperado:
\begin{itemize}
  \item Robustez: palabras y construcciones fuera de la gramática.
  \item Ambigüedad: cómo la maneja cada enfoque.
  \item Escalabilidad: costo de ampliar la gramática frente a reentrenar o cambiar el prompt.
  \item Costo: tiempo, recursos y dinero por oración.
\end{itemize}

% Responsable: P1  (versión completa en extendido/10_conclusiones.tex)
\section{Conclusiones}
\label{sec:conclusiones}

\begin{itemize}
  \item \textbf{Una CFG describe bien un fragmento del español, y solo ese fragmento.} Dentro de
    él, el parser es exacto en decisión, tipo de error y análisis; fuera de él, el 20\,\% del
    español correcto de nuestro conjunto queda rechazado sin análisis.
  \item \textbf{La concordancia es lo que aporta la CFG con rasgos.} Es el único sistema que
    rechaza las 13 oraciones con error de concordancia, dice por qué y no da falsas alarmas.
    Verificarla aparte, como en la fase semántica de un compilador, fue lo que la hizo manejable.
  \item \textbf{La ambigüedad es el límite común.} La gramática no es LL(1) por la frase
    preposicional; el parser acierta la unión en el 59\,\% de los casos, spaCy y el LLM en más,
    pero ninguno reconoce las oraciones ambiguas.
  \item \textbf{Los modelos estadísticos ganan en cobertura y pierden en garantías:} etiquetan
    casi todo, pero no deciden la gramaticalidad o la confunden con frecuencia.
  \item \textbf{Un LLM puede seguir reglas formales, pero no se puede confiar en que lo haga:}
    decide como las reglas (0.97--0.99), pero a veces lo justifica con árboles inválidos.
\end{itemize}

\paragraph{Limitaciones y trabajo futuro.} El conjunto tiene 73 oraciones etiquetadas por el
mismo equipo que escribió la gramática. El parser no retrocede, así que pierde algunas
oraciones con palabras ambiguas; un parser de gráfico resolvería eso y mostraría todas las
lecturas de las oraciones ambiguas. La gramática se puede ampliar con sujeto omitido y
coordinación, y el experimento con el LLM, repetirse con modelos más grandes.


% Referencias. Cada persona agrega las que cita; P1 unifica el estilo.
\begin{thebibliography}{9}

\bibitem{chomsky56}
N.~Chomsky.
\newblock Three models for the description of language.
\newblock \emph{IRE Transactions on Information Theory}, 2(3):113--124, 1956.

\bibitem{jurafsky}
D.~Jurafsky y J.~H. Martin.
\newblock \emph{Speech and Language Processing}.
\newblock Borrador de la 3.\textsuperscript{a} ed., disponible en línea.

\bibitem{hmu}
J.~E. Hopcroft, R.~Motwani y J.~D. Ullman.
\newblock \emph{Introduction to Automata Theory, Languages, and Computation}.
\newblock 3.\textsuperscript{a} ed. Addison-Wesley, 2006.

\bibitem{shieber}
S.~M. Shieber.
\newblock Evidence against the context-freeness of natural language.
\newblock \emph{Linguistics and Philosophy}, 8(3):333--343, 1985.

\end{thebibliography}


\end{document}
